\chapter{A worked bridge: orders, periodic complexes and spectra}
\label{ch:worked-bridge}

The preceding chapters provide several legitimate connections, each with a
specified domain of validity. We now place them beside one another in a single
calculation. Its purpose is to identify exactly which structures survive a
change of viewpoint. The example is elementary enough to calculate without
appealing to an existence theorem, but rich enough to distinguish a generic
semisimple algebra, an integral order, a singular special fiber, an infinite
projective resolution, and a positive finite-dimensional Hodge operator. It
also directly uses the nodal order occurring in the inherited research
program \citep{GRHRepo,SchreiberSurfaceI,SchreiberSurfaceII}.

\section{The order and its two complementary elements}

Let $R$ be a commutative integral domain, let $u\in R$ be nonzero, and set
\[
 \Lambda=R[X]/(X^2-uX),\qquad x=\overline X,\qquad y=u-x.
\]
Then $xy=yx=0$. The polynomial is monic, so every element of $\Lambda$ has a
unique expression $a+bx$, with $a,b\in R$. Relative to the basis $(1,x)$,
multiplication by $x$ and $y$ is represented by
\[
 M_x=\begin{pmatrix}0&0\\1&u\end{pmatrix},\qquad
 M_y=\begin{pmatrix}u&0\\-1&0\end{pmatrix}.
\]
These matrices exhibit the zero-composition condition as an identity over the
coefficient ring, before any choice of scalar field or inner product.

\begin{proposition}[An exact periodic complex]\label{prop:periodic-order}
There is an exact complex of free $\Lambda$-modules
\[
 \cdots\xrightarrow{x}\Lambda\xrightarrow{y}\Lambda
 \xrightarrow{x}\Lambda\xrightarrow{y}\Lambda\xrightarrow{x}\cdots.
\]
Here an arrow labelled by an element denotes multiplication by that element.
Moreover, with $M=\Lambda/x\Lambda$, the augmented sequence
\[
 \cdots\xrightarrow{x}\Lambda\xrightarrow{y}\Lambda
 \xrightarrow{x}\Lambda\longrightarrow M\longrightarrow0
\]
is a free resolution of $M$.
\end{proposition}
\begin{proof}
For $a+bx\in\Lambda$,
\[
 x(a+bx)=(a+ub)x,\qquad y(a+bx)=a(u-x)=ay.
\]
Uniqueness of the basis expansion gives
\[
 \ker M_x=R(x-u)=Ry,\qquad \im M_y=Ry.
\]
It also gives $\ker M_y=Rx$: the coefficient of $x$ in $ay$ is $-a$, so
$ay=0$ implies $a=0$. Finally $\im M_x=Rx$, since $x=x\cdot1$. Thus the
kernel of each arrow is the image of the preceding one. The quotient map
onto $M$ has kernel $x\Lambda=Rx$, which verifies exactness at the augmented
term. Every unaugmented term is the free module $\Lambda$ of rank one.
\end{proof}

The ring need not be a domain: indeed $xy=0$. This is not a defect in the
argument. The coefficient ring $R$ is a domain, whereas the order $\Lambda$
is a ring with zero divisors. Keeping those two statements separate prevents
an erroneous cancellation of $x$ or $y$. In particular, this is not a
regular-sequence Koszul resolution. It is a periodic resolution arising from
the explicit annihilator identities.

\section{Tensor products detect information that ranks omit}

The module $M$ is isomorphic to $R$ as an $R$-module, with $x$ acting as zero.
The action of $y$ on $M$ is consequently multiplication by $u$. Tensoring the
resolution with $M$ over $\Lambda$ produces the chain complex
\[
 \cdots\xrightarrow{0}R\xrightarrow{u}R
 \xrightarrow{0}R\xrightarrow{u}R\xrightarrow{0}R.
\]
The rightmost copy is in degree zero.

\begin{proposition}[An explicit derived obstruction]\label{prop:tor-node}
For $M=\Lambda/(x)$ as above,
\[
 \Tor^{\Lambda}_{2j+1}(M,M)\cong R/uR\quad(j\geq0),
 \qquad
 \Tor^{\Lambda}_{2j}(M,M)=0\quad(j\geq1).
\]
In degree zero $M\otimes_\Lambda M\cong R$.
\end{proposition}
\begin{proof}
The displayed free resolution computes the derived tensor product. Its odd
differentials are zero, and its positive even differentials are
multiplication by $u$. Since $R$ is a domain and $u\ne0$, the latter maps are
injective. The homology in positive even degree is therefore zero. In odd
degree the kernel is all of $R$ and the incoming image is $uR$, giving
$R/uR$. In degree zero the incoming differential is zero. This establishes
all the assertions directly from the definition of homology.
\end{proof}

If $u$ is a unit, the positive Tor groups vanish, and the generic splitting
is already available over $R$. If $u$ is a nonunit, the odd groups record a
defect supported on its vanishing locus. Their $R$-rank is zero, so taking
generic ranks erases this information completely. This is a concrete reason
to retain the exact category of lattices and derived functors when comparing
orders. A rational representation and its integral model can agree in rank
while disagreeing substantially in extension theory.

This example also guards against a misleading computational inference.
Efficient manipulation of the two small matrices computes this particular
resolution. It says nothing about a polynomial-time procedure for arbitrary
module classification, arbitrary syzygies, or a cryptographic search problem.
The presentation size, coefficient bit lengths, and demanded output are part
of any complexity statement.

\section{Generic splitting and singular specialization}

Over the fraction field $K$ of $R$, the two elements
\[
 e_x=x/u,\qquad e_y=y/u
\]
are orthogonal idempotents with $e_x+e_y=1$. Thus
$K\otimes_R\Lambda\cong K\times K$. The denominators matter: if $u$ is not
invertible in $R$, these idempotents need not belong to $\Lambda$ itself.
The generic decomposition therefore need not be an integral decomposition.
This is the same elementary obstruction that appears when an overorder
introduces idempotents unavailable in the smaller order.

For $R=k[u]$ or $k[[u]]$, reduction modulo $u$ gives
\[
 \Lambda/u\Lambda\cong k[X]/(X^2).
\]
The special fiber has a nonzero nilpotent, whereas the generic fiber is a
product of fields. The periodic maps reduce to multiplication by $X$ and
$-X$, whose image and kernel both equal $kX$. Thus exactness survives this
particular reduction. This conclusion follows from an explicit calculation;
it must not be generalized to a statement that every exact complex remains
exact after arbitrary nonflat base change. Proposition~\ref{prop:tor-node}
already demonstrates the contrary phenomenon for another tensor operation.

\section{The finite Hodge operator of the same complex}

We next work over $\R$ and evaluate the parameter at a real number $t$.
This is a finite matrix family. Equip every copy of $\R^2$ with the standard
Euclidean inner product, and take the consecutive differentials
\[
 A_t=\begin{pmatrix}0&0\\1&t\end{pmatrix},\qquad
 B_t=\begin{pmatrix}t&0\\-1&0\end{pmatrix}.
\]
Their composite vanishes for every $t$, and each matrix has rank one.
The middle Hodge Laplacian is
\[
 \Delta_t=A_tA_t^T+B_t^TB_t.
\]

\begin{proposition}[Exactness and an explicit metric gap]
For every $t\in\R$,
\[
 \Delta_t=(1+t^2)I_2.
\]
Its kernel is zero and its least eigenvalue is $1+t^2\geq1$.
For every $\eta>0$, replacing both differentials by
$\eta A_t,\eta B_t$ preserves their images, kernels and exactness but
replaces the Laplacian by $\eta^2(1+t^2)I_2$.
\end{proposition}
\begin{proof}
Direct multiplication gives
\[
 A_tA_t^T=\begin{pmatrix}0&0\\0&1+t^2\end{pmatrix},\qquad
 B_t^TB_t=\begin{pmatrix}1+t^2&0\\0&0\end{pmatrix}.
\]
Adding proves the identity. A scalar multiple of the identity has the stated
eigenvalues. Multiplication of a linear map by a nonzero scalar changes
neither its image nor its kernel. Its contribution to $AA^T+B^TB$ is
multiplied by the square of that scalar, which proves the final assertion.
\end{proof}

Several lessons are visible simultaneously. A singular algebraic fiber does
not force this chosen finite Hodge gap to close. Conversely, exactness and
the same algebraic rank stratum do not provide a positive lower bound when
the scale of the differential is allowed to vary. The inner product and
normalization of the operator are essential data. The latter statement is
the finite counterpart of the much larger distinction between a topological
index and a uniform physical spectral gap. It supplies a counterexample to
an inference; it does not model the interacting Yang--Mills Hamiltonian.

\section{How the cryptographic comparison enters}

The derived obstruction $R/uR$ and the cyclic hidden subgroup in the
cryptography chapter are both quotients of explicitly presented modules.
That common language is useful: one can write maps, calculate kernels and
compare exact sequences. Their computational roles are nevertheless
different. A presentation of $R/uR$ includes its relation $u$. In an elliptic
curve discrete logarithm instance $Q=dP$, the unknown $d$ is precisely the
parameter of the hidden kernel
\[
 H_d=\{(-dt,t):t\in\Z/q\Z\}.
\]
Writing $H_d$ symbolically does not compute $d$. Fourier sampling obtains
information about its annihilator using a quantum evaluation procedure for
the public map $(a,b)\mapsto aP+bQ$. A useful algorithmic bridge must explain
how representations of the relevant maps are obtained without first solving
the hidden problem.

Similarly, writing down a lattice basis after knowing a secret scalar is not
a reduction from an unknown scalar to lattice reduction. The basis of the
endomorphism relation lattice in the cryptography chapter is available when
the action of the specified endomorphism on the cyclic subgroup is known.
That availability is an input condition. The abstract existence of an
endomorphism in an order does not automatically provide its efficient
evaluation, its scalar action on a selected subgroup, or a short basis in a
specified metric.

\chapter{A quantitative synthesis and a research program}

\section{The information that each theory actually supplies}

The most useful common objects are modules with extra structure and complexes
of maps between them. The relevant extra structure changes from chapter to
chapter. An order-lattice is an integral module inside a rational
representation. A cryptographic lattice is supplied with a norm,
representation and probability distribution. An elliptic complex is supplied
with differential operators and an exact principal-symbol sequence away from
the zero section. A quantum theory is supplied with positive correlation
functionals, spacetime covariance, locality and a reconstruction procedure.

These additions are not optional embellishments to an otherwise complete
algebraic answer. They determine which questions are meaningful. The matrix
equation $BA=0$ determines a variety of complexes, but no quantum resource
cost. A $K$-theory class determines the stable topological input to an index
pairing, but no absolute time scale. A free abelian group determines neither
the length of a shortest vector nor the distribution of a sampled error.

\begin{longtable}{p{.20\textwidth}p{.33\textwidth}p{.37\textwidth}}
\toprule
Interface & Data retained & Additional data needed for the next question\\
\midrule\endhead
Order to rational algebra & Scalar extension, representation, rank & Integral
conductor, local structure and lattice class\\
Order to geometric lattice & An embedding into a real vector space & Positive
form, basis, coefficient sizes and distortion bounds\\
Complex to cohomology & Kernel modulo image & Chosen inner products and
quantitative singular values\\
Elliptic symbol to index & Symbol $K$-class and characteristic pairing & Domains,
boundary conditions and estimates on positive spectrum\\
Group action to quantum algorithm & An abstract operation or oracle & Reversible
implementation, cleanup, precision and total resource accounting\\
Lattice action to LWE problem & Module arithmetic & Error law, parameter regime,
sample distribution and a reduction\\
Euclidean correlations to a QFT & Positive continuum functionals and axioms &
Control of phases, quantitative large-distance decay and spectral support\\
\bottomrule
\end{longtable}

The table organizes a rigorous synthesis because each row is an interface
with testable requirements. It prevents a conclusion from one category being
silently reinterpreted as a conclusion in another.

\section{A general comparison estimate}

There is a simple quantitative principle behind several of these interfaces.
Let $V$ be a finite-dimensional real inner-product space, let $\Delta$ be a
positive semidefinite self-adjoint operator, and suppose its kernel is $K$.
A gap statement is a lower bound on the orthogonal complement of $K$:
\[
 \langle v,\Delta v\rangle\geq c\|v\|^2\qquad(v\perp K).
\]
Algebraic identifications transport vectors. To transport this inequality,
they must also control norms and quadratic forms.

\begin{proposition}[Transport of a coercivity bound]
Let $V,W$ be finite-dimensional real inner-product spaces, $T:W\to V$ a
linear isomorphism, and $q$ a nonnegative quadratic form on $V$ with kernel
$K$. Suppose
\[
 q(v)\geq c\,\operatorname{dist}(v,K)^2\quad(v\in V),
 \qquad \|Tw\|\geq m\|w\|\quad(w\in W)
\]
with $c,m>0$. Put $q_T(w)=q(Tw)$ and $K_T=T^{-1}K$. Then
\[
 q_T(w)\geq cm^2\,\operatorname{dist}(w,K_T)^2.
\]
\end{proposition}
\begin{proof}
For every $k\in K$, write $k=Tz$ with $z\in K_T$. Then
$\|Tw-k\|=\|T(w-z)\|\geq m\|w-z\|$. Taking the infimum over $k$, equivalently
over $z$, gives
$\operatorname{dist}(Tw,K)\geq m\operatorname{dist}(w,K_T)$.
Apply the assumed bound for $q$ and square the latter inequality.
\end{proof}

This estimate has deliberately modest scope. It does not construct $T$, prove
that $m$ stays away from zero in a limit, or identify a finite matrix with a
physical Hamiltonian. It shows exactly what must be added to an algebraic
identification to transport a lower spectral bound. The same smallest
singular value controls distortion of short vectors in a lattice. The same
absence of a uniform bound can invalidate a proposed continuum-limit argument
or a proposed cryptographic reduction.

\begin{corollary}[Why pointwise equivalence is insufficient]
For a family of isomorphisms $T_j$, positive constants $m_j$ for each fixed
$j$ do not imply a uniform positive lower bound for the transported quadratic
forms. On $V=W=\mathbb R$, with $j\ge1$ and $q(v)=\|v\|^2$,
the choice $T_j=j^{-1}I$ gives
$q_{T_j}(w)=j^{-2}\|w\|^2$.
\end{corollary}
\begin{proof}
The displayed identity follows from homogeneity of the norm. Its least
Rayleigh quotient is $j^{-2}$, which tends to zero. Each $T_j$ remains an
isomorphism.
\end{proof}

\section{A diagram with explicit logical boundaries}

There is a defensible route from the two books to the other topics. Integral
orders organize endomorphism rings and modules. These produce complexes and
deformation problems, whose finite algebra can be studied through rank
strata. Smooth geometric counterparts have symbol complexes and index
classes. Gauge theory supplies a particularly important elliptic deformation
complex. Finite abelian groups and modules also support the character theory
underlying quantum hidden-subgroup algorithms. The links are structural and,
in examples such as the preceding chapter, computationally explicit.

What is absent is a theorem making the path transitive at the level of the
most difficult conclusions. An endomorphism order does not turn an ASD index
into an ECC discrete logarithm algorithm. A positive finite Hodge Laplacian
does not establish a uniform infinite-volume Yang--Mills gap. A functional
equation and real coefficients do not identify an arithmetic determinant
with a self-adjoint operator. These omissions are the mathematical content of
the remaining research questions, rather than details to be supplied by an
unlabelled analogy.

\section{Three concrete research directions}

\subsection{Arithmetic modules with explicit metric control}

A productive first direction is to choose a family of orders and record,
together, its integral presentation, archimedean embeddings, positive form,
local conductor, and a basis with certified coefficient sizes. The codifferent
calculation in the cryptography chapter provides the correct dual object
under a trace pairing. The shear example shows why a module isomorphism alone
cannot certify a short-vector comparison. A useful theorem would provide a
uniform, explicit distortion bound for a specified family of transformations,
with the bound expressed in the same parameter as the input bit length.

For an imaginary quadratic order, one can begin with explicit ideal lattices
and the complex absolute-value metric. One must state which ideals are
invertible, whether the order is maximal, how the discriminant varies, and
which basis or reduced representative is produced. These restrictions are
essential for a valid theorem about dual ideals or norm comparison. Extending
such a theorem to noncommutative orders requires specifying left versus right
modules and the chosen involution or positive representation.

The computational companion provides small exact examples, not an asymptotic
complexity result. A subsequent study should measure both coefficient growth
and geometric distortion. Reporting only a favorable rank or module
decomposition would leave the main quantitative issue unanswered.

\subsection{Resource-certified elliptic-curve quantum arithmetic}

The reproducibility program is implemented in the companion's
\texttt{workbench/program3/ecc.py}, with separate emitted variable-point
addition oracles, independent gate counting and clean-work verification.
The current continuation, \emph{Arithmetic, Coverings and Learned Quantum
Renormalization}, gives the construction, proofs, tested resource comparisons
and observational controls. Historical editions retain the original research
direction; this source records its implemented successor.

The comparison fixes nonsingular finite curves, affine coordinates with a
unique infinity encoding, invalid-word behavior, ideal X/CX/CCX gates and
ancillary-register initialization. Both the binary path and coherent signed
endomorphism decomposition include every preparation and cleanup operation.
Finite exhaustive-domain checks and actual cuStateVec replay are recorded at
their precise scope. The bounded truth-table compiler is not a polynomial-size
implementation for cryptographic curves, and logical counts are not
error-corrected physical resource estimates.

In the tested instances, decomposition does not reduce the complete budget.
The emitted circuits make that result reproducible instead of inferring a
saving from shorter scalar components alone. The independent source, circuit
streams, assumption manifest and verification reports accompany the versioned
continuation. Its hierarchical channel and covering-theory extensions preserve
the distinction between an arithmetic endomorphism, a quantum channel, a graph
cover and a learned renormalization map.

\subsection{From elliptic estimates to a uniform physical gap}

A third direction begins with the Yang--Mills existence framework reviewed in
this manuscript. The independent assessment separates existence from the
mass-gap problem \citep{YMReview2026}. A further theorem needs a quantitative
bound strong enough to survive ultraviolet and infinite-volume limiting
procedures, together with a correct identification of the reconstructed
observables and vacuum sector.

The spectral theorem argument in the gauge chapter gives a precise target:
uniform exponential decay on the orthogonal complement of the vacuum is
equivalent to a positive lower spectral threshold, with the stated domain
and cyclicity conditions. The UV counterexample there proves that matching
leading short-distance behavior cannot alone establish this target. A new
mass-gap argument must therefore add infrared information. The ASD index
formula may organize classical deformation spaces, but it does not provide
the required estimate for the interacting quantum Hamiltonian.

Formalization could assist by making the limiting and reconstruction
interfaces explicit. The hard step would remain a proved analytic estimate,
not a declaration naming a gap constant. A sound Lean statement must quantify
over the relevant cutoffs and volumes, give the correct normalization, and
make positivity of the common lower bound an actual conclusion or a clearly
identified hypothesis. The formalization appendix explains this distinction
using smaller declarations that can be compiled in the present project.

\section{What this monograph establishes}

This work supplies a coherent collection of explicit constructions,
elementary proofs, counterexamples to invalid transfers, a source audit,
selected machine-checked algebra, and reproducible finite experiments. It
does not assert a proof of GRH, a Yang--Mills mass gap, a generally faster ECC
algorithm, or a new security reduction for lattice cryptography. Those would
require additional theorems not contained in the source material or proved
here.

The strongest common principle obtained from the comparison is quantitative:
an algebraic or categorical connection becomes an algorithmic or spectral
connection only after its representations, metrics, probability laws,
normalizations and limiting operations are controlled. The examples make
each part of that principle testable. The research directions above specify
what a successful next theorem would have to say and what evidence would
make it independently assessable.
